% =========================================================
% Paste-ready blocks for chapters/05_results.tex, Section 5.2.2 (your numbering, your text kept as is).
% Block 3 is the paragraph that summarizes the full-posterior bar chart, which is in the appendix
% (appendix_B_posterior_comparison.tex: result_5_2/fig_5_8_posterior_shift.pdf, label fig:posterior_shift, plus
% Table tab:candidate_pool and the detailed description). Block 3 replaces the figure environment and the three
% paragraphs after Table 5.2 and before "Score Margins of the Constructed Pairs".
% Blocks 1, 2 and 4 are optional: 1 completes the unfinished "Second, ..." sentence of the 5.2.2 intro,
% 2 replaces the leftover sentence fragments after Table 5.2 (the table now shows pair j = 2, so the
% scores are 0.83 / 0.24 and 1.48 / 0.31). Block 4 is the closing paragraph that answers both parts of RQ1
% and names the two qualifications (classifier's Telling performance, move distribution of the source data).
% Block 5 points to the appendix file (figure, table of the candidate pool, full description).
% Everything else (seven-context figure, model composition, utility swap) goes to the appendix and is
% not referenced here. Replace the label names by yours if they differ.
% =========================================================

% ---------- 1. completion of the 5.2.2 intro ("Second, ...") ----------
In this section, we first illustrate how the preference pairs are constructed with the error-conditioned tutor-move scoring of Sections~\ref{subsec:error_conditioned_preference_scoring} and~\ref{subsec:preference_pair_construction}, using two contexts with different student error types (Figure~\ref{fig:context_ranking} and Table~\ref{tab:example_pairs}). Second, we compare the move distributions of the chosen and the rejected responses over all 24,298 pairs (Figure~\ref{fig:posterior_shift} in Appendix~\ref{app:posterior_comparison}). Third, we report the score margins that separate the two sides of the pairs (Figure~\ref{fig:margin_by_pair}).

% ---------- 2. sentence after Table 5.2 (replaces the leftover fragments; scores of pair j = 2) ----------
Table~\ref{tab:example_pairs} lists the second pair (\(j = 2\)) of both contexts. In both pairs, one response asks the student which quantity to recompute and the other gives the complete calculation. The scoring puts the question on the chosen side and the complete calculation on the rejected side for the \textit{numerical calculation error} (scores 0.83 against 0.24), and the complete calculation on the chosen side and the question on the rejected side for \textit{missing or incorrect factual knowledge} (1.48 against 0.31). The same two kinds of response thus receive opposite preferences under the two error types, which is the intended behavior of the utility \(u(m,e)\).

% ---------- 3. the paragraph that replaces the bar chart in 5.2.2 (the figure itself is in the appendix; no figure environment here) ----------
\paragraph{Chosen--Rejected Probability Comparison}
Figure~\ref{fig:posterior_shift} in Appendix~\ref{app:posterior_comparison} extends the comparison from the two contexts to all 24,298 pairs. For each error type and each move, it compares the mean classifier posterior of the chosen responses with that of the rejected responses. In all six error types with a \textit{Focus} or \textit{Telling} utility, the differences follow the utility table: the move with the largest utility is higher on the chosen side, and no move with zero utility is. For the three \textit{Focus}-associated error types, the mean \textit{Focus} posterior is 0.88--0.90 on the chosen side against 0.68--0.76 on the rejected side. For the three \textit{Telling}-associated error types, \textit{Telling} is higher on the chosen side (0.11--0.13 against 0.02), but the largest utility is not reflected in the largest posterior mass: \textit{Focus} remains the most probable move of the chosen responses (0.69--0.78). The utility table is thus reflected in the difference between the two sides, not in the dominant move of the chosen responses. \textit{Misunderstanding of the question} is the exception: its only positive utility belongs to \textit{Generic}, which the candidates almost never realize, so the two sides differ only marginally (0.004 against 0.002) and the scoring has nothing to prefer.

Two properties of the data limit the \textit{Telling} side. First, the candidate pool is nearly the same for every error type and dominated by \textit{Focus}, the top label of 45,104 of the 48,839 candidates; \textit{Telling} is the top label of 1,498, and 307 of the 693 contexts hold no \textit{Telling}-labelled candidate at all (Table~\ref{tab:candidate_pool} in Appendix~\ref{app:posterior_comparison}). Second, the moves of the candidates are those assigned by the classifier of Section~\ref{subsec:selected_classifier}, which recovers about half of the \textit{Telling} utterances of its test set (F1 of 0.40), and the \textit{Telling} utilities themselves rest on one to four first tutor turns (Figure~\ref{fig:associations}). Both points are discussed in Chapter~\ref{chapter:Discussions}.

% ---------- 4. closing paragraph (after "Score Margins of the Constructed Pairs") ----------
\paragraph{The Preference Signal in the Constructed Dataset}
Taken together, the results answer Research Question~1. As to the associations, the first tutor turns of the 693 training contexts are dominated by \textit{Generic} for every error type, and it is only after the base rate of each move is divided out that the error types differ: PPMI associates \textit{Focus} with \textit{numerical calculation error}, \textit{none of the above}, and \textit{reached the correct solution but proceeded further}, \textit{Telling} with \textit{missing or extra quantity}, \textit{missing or incorrect factual knowledge}, and \textit{unit conversion error}, and \textit{Generic} alone with \textit{misunderstanding of the question} (Section~\ref{subsec:observed_associations}). As to their reflection in the constructed dataset, the same kind of response, a complete calculation, is rejected for a \textit{numerical calculation error} and chosen for \textit{missing or incorrect factual knowledge} (Figure~\ref{fig:context_ranking} and Table~\ref{tab:example_pairs}); over all 24,298 pairs, the mean posteriors of the chosen and the rejected responses differ in the direction of the utility table in every error type (Figure~\ref{fig:posterior_shift} in Appendix~\ref{app:posterior_comparison}); and the separation between the two sides is measurable, with a mean normalized margin of the first pair between 0.37 and 0.71 in the six error types with a \textit{Focus} or \textit{Telling} utility (Figure~\ref{fig:margin_by_pair}). The constructed dataset therefore carries an error-conditioned preference signal for six of the seven error types; for \textit{misunderstanding of the question}, whose only positive utility belongs to a move that the candidates do not realize, it does not. What the scoring transfers to the dataset is primarily the sign pattern of the association, that is, which moves occur above their base rate for an error type; the signal is a difference in posterior mass while \textit{Focus} remains the most probable move of most responses on both sides, and for the \textit{Telling}-associated error types it is bounded by the candidate pool, rests on few observations, and is measured with a classifier that recovers about half of the \textit{Telling} utterances. These properties of the source data, of the candidate pool, and of the classifier are discussed in Chapter~\ref{chapter:Discussions}.

% ---------- 5. appendix: figure fig:posterior_shift + table tab:candidate_pool + full description ----------
% See appendix_B_posterior_comparison.tex (a \section for pages/appendix.tex, after the section of appendix_A_contexts.tex).
% Block 3 above cites both by their labels (fig:posterior_shift, tab:candidate_pool, app:posterior_comparison).

% ---------- 6. "Score Margins of the Constructed Pairs": paragraph and figure environment for the new fig_5_11 (mean, first 30 pairs) ----------
\paragraph{Score Margins of the Constructed Pairs}
Figure~\ref{fig:margin_by_pair} shows the score margins that the pairing rule of Section~\ref{subsec:preference_pair_construction} produces. Because pair \(j\) joins rank \(j\) with rank \(n-j+1\), the margin decreases with \(j\) in every context. The figure reports the mean normalized margin \(\Delta_{e,j}/\max_{m}\operatorname{PPMI}(m,e)\) over all contexts of each error type for the first 30 pairs of a context, the number of pairs that every context has. For the first pair, the mean lies between 0.37 (\textit{missing or incorrect factual knowledge}) and 0.71 (\textit{numerical calculation error} and \textit{none of the above}) in six of the seven error types; it is 0.03 for \textit{misunderstanding of the question}, where the three other utilities are zero and the scores of all candidates are bounded by the \textit{Generic} utility of 0.07 (third row of Figure~\ref{fig:app_contexts}). By the tenth pair the mean has fallen to 0.11--0.28, and by the thirtieth pair to 0.04 or less in every error type. The two error types with the largest \textit{Focus} utilities (\textit{numerical calculation error}, 0.90; \textit{none of the above}, 0.30) keep the widest margins at every \(j\); the other four error types with a \textit{Focus} or \textit{Telling} utility lie between 0.37 and 0.49 at the first pair and between 0.11 and 0.17 at the tenth.

    \begin{figure}[htbp]
        \centering
        \includegraphics[width=\textwidth]{figures/fig_5_11_margin_by_pair.pdf}
        \caption[Score margins of the constructed pairs by pair index.]{Score margin of the constructed pairs by pair index \(j\) within a context (pair \(j\) joins rank \(j\) with rank \(n-j+1\)). Each point is the mean normalized margin \(\Delta_{e,j}/\max_{m}\operatorname{PPMI}(m,e)\) over all contexts of the error type. Only the first 30 pairs of every context are shown, so that every point rests on all contexts of its error type (87, 60, 49, 167, 97, 34, and 199 contexts in the order of the legend).}
        \label{fig:margin_by_pair}
    \end{figure}
